% ============================================================
%  Laporan Tugas Kecil 2 IF2211 Strategi Algoritma
%  Voxelization Objek 3D menggunakan Octree
%  Zahran Alvan Putra Winarko - 13524124
% ============================================================

\documentclass[12pt, a4paper]{report}

% ── Encoding & Bahasa ────────────────────────────────────────
\usepackage[T1]{fontenc}
\usepackage[utf8]{inputenc}
\usepackage[bahasa]{babel}

% ── Tata Letak Halaman ───────────────────────────────────────
\usepackage[
  top=3cm, bottom=2.5cm,
  left=3cm, right=2.5cm
]{geometry}
\usepackage{setspace}
\onehalfspacing

% ── Font & Tipografi ─────────────────────────────────────────
\usepackage{lmodern}
\usepackage{microtype}
\usepackage{parskip}
\setlength{\parskip}{6pt}

% ── Warna & Grafik ───────────────────────────────────────────
\usepackage[dvipsnames, table]{xcolor}
\usepackage{graphicx}
\usepackage{float}
\usepackage{tikz}

% Warna kustom
\definecolor{itbblue}{RGB}{0, 70, 127}
\definecolor{codeblue}{RGB}{0, 100, 180}
\definecolor{codegray}{RGB}{245, 245, 245}
\definecolor{codecomment}{RGB}{100, 100, 100}
\definecolor{codestring}{RGB}{163, 21, 21}
\definecolor{codekeyword}{RGB}{0, 0, 255}

% ── Listing (Kode Program & Pseudocode) ──────────────────────
\usepackage{listings}

\lstdefinestyle{pseudocode}{
  backgroundcolor=\color{codegray},
  basicstyle=\ttfamily\small,
  keywordstyle=\bfseries,
  commentstyle=\itshape\color{codecomment},
  numbers=left,
  numberstyle=\tiny\color{gray},
  numbersep=8pt,
  frame=single,
  framerule=0.5pt,
  rulecolor=\color{itbblue!40},
  breaklines=true,
  captionpos=b,
  tabsize=4,
  showstringspaces=false,
  morekeywords={FUNGSI, JIKA, MAKA, UNTUK, DARI, HINGGA, RETURN,
                JIKA, TIDAK, ADA, LAIN, TULIS, BACA, DAN, ATAU,
                DIVIDE, CONQUER, PARALEL, TUNGGU, SEMUA, SELESAI,
                SEQUENTIAL, return, true, false, and, or, not,
                function, if, then, while, do, else},
}

\lstdefinestyle{cpp}{
  language=C++,
  backgroundcolor=\color{codegray},
  basicstyle=\ttfamily\small,
  keywordstyle=\color{codekeyword}\bfseries,
  commentstyle=\itshape\color{codecomment},
  stringstyle=\color{codestring},
  numbers=left,
  numberstyle=\tiny\color{gray},
  numbersep=8pt,
  frame=single,
  framerule=0.5pt,
  rulecolor=\color{itbblue!40},
  breaklines=true,
  captionpos=b,
  tabsize=4,
  showstringspaces=false,
}

% ── Tabel ────────────────────────────────────────────────────
\usepackage{booktabs}
\usepackage{tabularx}
\usepackage{multirow}
\usepackage{array}
\usepackage{longtable}

% ── Header & Footer ──────────────────────────────────────────
\usepackage{fancyhdr}
\pagestyle{fancy}
\fancyhf{}
\fancyhead[L]{\small\textit{Laporan IF2211 Strategi Algoritma}}
\fancyhead[R]{\small\textit{Voxelization 3D menggunakan Octree}}
\fancyfoot[C]{\small Laporan IF2211 Strategi Algoritma - Teknik Informatika ITB
              - Semester 2 Tahun 2025/2026 \hfill \thepage}
\renewcommand{\headrulewidth}{0.4pt}
\renewcommand{\footrulewidth}{0.4pt}

% Halaman chapter (plain style) juga punya footer
\fancypagestyle{plain}{
  \fancyhf{}
  \fancyfoot[C]{\small Laporan IF2211 Strategi Algoritma - Teknik Informatika ITB
                - Semester 2 Tahun 2025/2026 \hfill \thepage}
  \renewcommand{\headrulewidth}{0pt}
  \renewcommand{\footrulewidth}{0.4pt}
}

% ── Judul Bab & Bagian ───────────────────────────────────────
\usepackage{titlesec}
\titleformat{\chapter}[display]
  {\normalfont\Large\bfseries\color{itbblue}}
  {BAB \thechapter}{0pt}{\Large}
\titlespacing*{\chapter}{0pt}{0pt}{20pt}

\titleformat{\section}
  {\normalfont\large\bfseries\color{itbblue}}
  {\thesection}{1em}{}

\titleformat{\subsection}
  {\normalfont\normalsize\bfseries}
  {\thesubsection}{1em}{}

% ── Daftar Isi ───────────────────────────────────────────────
\usepackage{tocloft}
\renewcommand{\cftchapfont}{\bfseries}
\renewcommand{\cftchappagefont}{\bfseries}
\setlength{\cftbeforechapskip}{4pt}

% ── Hyperlink ────────────────────────────────────────────────
\usepackage[
  colorlinks=true,
  linkcolor=itbblue,
  urlcolor=itbblue,
  citecolor=itbblue,
  pdftitle={Laporan Tucil 2 IF2211 - Voxelization 3D},
  pdfauthor={Zahran Alvan Putra Winarko}
]{hyperref}

% ── Algoritma (Notasi Algoritmik) ────────────────────────────
\usepackage{algorithm}
\usepackage{algpseudocode}
\algrenewcommand\algorithmicprocedure{\textbf{Fungsi}}
\algrenewcommand\algorithmicif{\textbf{Jika}}
\algrenewcommand\algorithmicelse{\textbf{Lain}}
\algrenewcommand\algorithmicthen{\textbf{maka}}
\algrenewcommand\algorithmicfor{\textbf{Untuk}}
\algrenewcommand\algorithmicdo{\textbf{lakukan}}
\algrenewcommand\algorithmicreturn{\textbf{return}}
\algrenewcommand\algorithmicwhile{\textbf{Selama}}
\algrenewcommand\algorithmicend{\textbf{Akhir}}

% ── Misc ─────────────────────────────────────────────────────
\usepackage{enumitem}
\usepackage{amsmath}
\usepackage{caption}
\usepackage{subcaption}
\usepackage{pifont}   % \ding{51} = centang
\usepackage{mdframed}

\newcommand{\cmark}{\ding{51}}
\newcommand{\xmark}{\ding{55}}

% ── Info Box ─────────────────────────────────────────────────
\newmdenv[
  linecolor=itbblue!60,
  linewidth=1pt,
  backgroundcolor=itbblue!5,
  innerleftmargin=10pt,
  innerrightmargin=10pt,
  innertopmargin=8pt,
  innerbottommargin=8pt,
  skipabove=6pt,
  skipbelow=6pt,
  splittopskip=\baselineskip,
  splitbottomskip=\baselineskip,
]{infobox}

% ─────────────────────────────────────────────────────────────
%  DOKUMEN MULAI
% ─────────────────────────────────────────────────────────────
\begin{document}

% ── Halaman Sampul ───────────────────────────────────────────
\begin{titlepage}
  \centering
  \vspace*{1cm}

  {\LARGE\bfseries LAPORAN TUGAS KECIL 2\\
   IF2211 STRATEGI ALGORITMA\par}

  \vspace{1.2cm}

  {\Large\bfseries
   VOXELIZATION OBJEK 3D\\
   MENGGUNAKAN OCTREE\par}

  \vspace{1.5cm}
  % Logo ITB — ganti path jika perlu, atau hapus baris ini
  % \includegraphics[width=4cm]{logo_itb.png}
  % \vspace{1cm}

  \vspace{0.8cm}
  Disusun oleh:

  \vspace{0.4cm}
  \begin{tabular}{c}
    \textbf{Yusuf Faishal Listyardi}\\
    NIM: 13524014\\[0.5em]
    \textbf{Zahran Alvan Putra Winarko}\\
    NIM: 13524124
  \end{tabular}

  \vfill

  {\bfseries PROGRAM STUDI TEKNIK INFORMATIKA}\\
  SEKOLAH TEKNIK ELEKTRO DAN INFORMATIKA\\
  INSTITUT TEKNOLOGI BANDUNG\\
  2026
\end{titlepage}

% ── Daftar Isi / Tabel / Gambar ──────────────────────────────
\tableofcontents
\listoftables
\listoffigures

% =============================================================
%  BAB 1 – PENDAHULUAN
% =============================================================
\chapter{Dekripsi Masalah}

Representasi tiga dimensi digital pada komputer dapat dilakukan dalam berbagai
format. Salah satu format yang paling umum adalah \textit{polygon mesh}, yaitu
kumpulan segitiga yang membentuk permukaan objek, umumnya disimpan dalam
format \texttt{.obj}. Di sisi lain, representasi berbasis \textit{voxel} —
elemen volumetrik berbentuk kubus seperti piksel dalam ruang tiga dimensi —
semakin banyak digunakan dalam berbagai aplikasi, mulai dari simulasi fisika,
pencetakan 3D, visualisasi medis, hingga permainan video bergaya \textit{blocky}
seperti Minecraft.

\begin{figure}[H]
  \centering
  \includegraphics[width=0.7\textwidth]{voxel.png}
  % \fbox{\parbox{0.85\textwidth}{\centering\vspace{2.5cm}
  %   [Screenshot terminal: perintah + output statistik cube depth 2]
  %   \vspace{2.5cm}}}
  \caption{Visualisasi voxelization}
\end{figure}

\textit{Voxelization} adalah proses mengonversi objek 3D kontinu (polygon mesh)
menjadi representasi diskrit berbasis voxel. Tantangan utamanya adalah efisiensi:
sebuah model bisa memiliki ratusan ribu segitiga, dan menguji setiap segitiga
terhadap setiap calon voxel secara langsung memiliki kompleksitas
$O(F \times N^3)$ di mana $F$ adalah jumlah \textit{face} dan $N$ adalah
resolusi tiap dimensi, yang sangat tidak praktis untuk ukuran besar.

Solusi efisien untuk masalah ini adalah menggunakan struktur data \textbf{Octree}
— sebuah pohon di mana setiap node internal memiliki tepat 8 anak — dikombinasikan
dengan paradigma \textbf{Divide and Conquer}. Dengan Octree, ruang tiga dimensi
dibagi secara rekursif menjadi delapan oktant hingga resolusi yang diinginkan
tercapai, dan hanya cabang yang bersinggungan dengan permukaan objek yang
ditelusuri lebih lanjut.



% =============================================================
%  BAB 2 – ALGORITMA
% =============================================================
\chapter{ALGORITMA}

\section{Deskripsi Umum Program}

Algoritma yang digunakan adalah \textbf{Divide and Conquer} berbasis
\textbf{Octree}. Ide dasarnya sederhana: daripada memeriksa setiap segitiga
terhadap setiap voxel potensial, kita memulai dari satu kotak besar yang
mencakup seluruh model, kemudian membaginya secara rekursif menjadi delapan
oktant yang lebih kecil. Pada setiap pembagian, kita hanya meneruskan
segitiga-segitiga yang benar-benar berpotongan dengan oktant tersebut ke
level berikutnya. Oktant yang tidak berisi segitiga langsung dibuang
(\textit{pruned}), sehingga komputasi terkonsentrasi hanya pada bagian ruang
yang benar-benar terisi permukaan objek.

\section{Struktur Data Octree}

Octree adalah pohon (\textit{tree}) di mana setiap node internal memiliki
tepat \textbf{8 anak}. Setiap node merepresentasikan sebuah
\textit{Axis-Aligned Bounding Box} (AABB) — kotak yang sisi-sisinya sejajar
dengan sumbu koordinat X, Y, dan Z.

\begin{infobox}
\textbf{Analogi:} Bayangkan sebuah kubus besar. Jika kita belah dengan tiga
bidang tegak lurus (masing-masing sejajar dengan sumbu X, Y, dan Z tepat di
tengah), kita mendapatkan 8 kubus kecil berukuran setengahnya. Inilah satu
langkah subdivisi Octree.
\end{infobox}

\subsection{Inisialisasi Root}

Root AABB dihitung sebagai bounding box seluruh vertex model, kemudian
dibuat menjadi \textbf{kubus sempurna} menggunakan sisi terpanjang dari
ketiga dimensi (ditambah faktor \textit{padding} 0,1\%). Hal ini menjamin
bahwa semua voxel yang dihasilkan berukuran seragam (\textit{uniform voxels}).

\section{Algoritma Divide and Conquer}

Algoritma bekerja dalam tiga tahap khas Divide and Conquer:

\begin{description}[leftmargin=2.5cm, labelwidth=2cm]
\item[\textbf{DIVIDE}]
    Program terlebih dahulu membentuk \textit{root bounding box} yang
    mencakup seluruh model 3D, lalu menyesuaikannya menjadi kubus agar
    ukuran voxel seragam. Setelah itu, pada setiap pemanggilan
    \texttt{buildRecursive()}, ruang AABB saat ini dibagi menjadi 8
    sub-ruang yang lebih kecil menggunakan fungsi \texttt{subdivide()}.
    Karena tiap sub-ruang merupakan versi lebih kecil dari persoalan
    semula, maka proses voxelisasi dapat dipecah menjadi delapan
    upa-persoalan yang sejenis.

\item[\textbf{CONQUER}]
    Masing-masing upa-persoalan diselesaikan dengan memeriksa segitiga
    mana saja yang masih berpotongan dengan setiap oktan anak menggunakan
    uji \textit{Triangle-AABB intersection} berbasis
    \textit{Separating Axis Theorem} (SAT). Jika suatu oktan tidak lagi
    memiliki segitiga aktif, rekursi pada cabang tersebut dihentikan
    (\textit{pruning}). Jika kedalaman maksimum telah tercapai, AABB
    tersebut langsung dinyatakan sebagai sebuah voxel. Jika belum,
    prosedur yang sama dipanggil kembali secara rekursif pada setiap
    oktan anak. Pada kedalaman awal, penyelesaian delapan sub-masalah ini
    dapat dijalankan paralel menggunakan \texttt{std::async} agar waktu
    komputasi lebih cepat.

\item[\textbf{COMBINE}]
    Solusi dari tiap sub-persoalan digabungkan dengan cara mengumpulkan
    seluruh \textit{leaf node} valid hasil rekursi menjadi himpunan voxel
    akhir program. Setiap cabang yang mencapai kondisi \textit{leaf}
    menambahkan AABB-nya ke dalam container \texttt{voxels}, sehingga
    gabungan semua hasil anak membentuk representasi voxel dari model
    asli. Pada bagian paralel, seluruh \textit{thread} anak terlebih
    dahulu disinkronkan dengan \texttt{future::get()}, sehingga hasil
    dari kedelapan cabang benar-benar telah terkumpul sebelum proses
    kembali ke level induknya.
\end{description}

\textbf{Kondisi basis:}
\begin{itemize}
  \item Jika daftar \textit{face} kosong: node di-\textit{prune},
        dicatat sebagai node yang tidak perlu ditelusuri.
  \item Jika \texttt{depth == maxDepth}: node dicatat sebagai
        satu \textbf{voxel} output.
\end{itemize}

\section{Pseudocode}

Berikut adalah pseudocode fungsi inti algoritma:

\begin{lstlisting}[style=pseudocode,
  caption={Pseudocode Algoritma Voxelization Divide and Conquer},
  label={lst:pseudo}]
FUNGSI build():
    Hitung AABB root dari seluruh vertex model
    Jadikan root berbentuk kubus sempurna (padding 0.1%)
    allFaces <- [0, 1, 2, ..., F-1]
    buildRecursive(rootAABB, allFaces, depth = 0)

FUNGSI buildRecursive(box, faces, depth):
    // Catat statistik node yang terbentuk
    JIKA depth > 0 MAKA
        nodesFormed[depth]++

    // BASIS 1: Tidak ada permukaan pada area ini
    JIKA faces KOSONG MAKA
        JIKA depth > 0 MAKA
            nodesSkipped[depth]++
        RETURN

    // BASIS 2: Kedalaman maksimum tercapai
    JIKA depth == maxDepth MAKA
        tambahkan box ke daftar voxels
        voxelCount++
        RETURN

    // DIVIDE: Bagi box saat ini menjadi 8 upa-persoalan
    subBoxes[8] <- box.subdivide()
    childFaces[8] <- array of list kosong

    // Persiapkan data untuk tiap sub-persoalan
    UNTUK setiap faceIdx dalam faces:
        v0, v1, v2 <- vertex dari face[faceIdx]
        UNTUK i dari 0 hingga 7:
            JIKA subBoxes[i].intersectsTriangle(v0, v1, v2) MAKA
                tambahkan faceIdx ke childFaces[i]

    // CONQUER: Selesaikan masing-masing upa-persoalan
    JIKA depth < PARALLEL_DEPTH_LIMIT MAKA
        UNTUK i dari 0 hingga 7:
            threads[i] <- async(buildRecursive(subBoxes[i],
                                               childFaces[i],
                                               depth + 1))
        UNTUK i dari 0 hingga 7:
            tunggu threads[i] selesai
    JIKA TIDAK MAKA
        UNTUK i dari 0 hingga 7:
            buildRecursive(subBoxes[i], childFaces[i], depth + 1)

    // COMBINE: Hasil dari seluruh anak tergabung dalam daftar voxels
    // karena setiap leaf valid menambahkan box-nya ke container global
    RETURN

AKHIR FUNGSI

\end{lstlisting}

\section{Uji Interseksi Triangle-AABB: Separating Axis Theorem}

Untuk menentukan apakah sebuah segitiga berpotongan dengan sebuah
\textit{Axis-Aligned Bounding Box} (AABB), program ini menggunakan
\textit{Separating Axis Theorem} (SAT). Teorema ini menyatakan bahwa dua
buah bangun konveks \textbf{tidak berpotongan jika dan hanya jika}
terdapat minimal satu sumbu di mana proyeksi kedua bangun \textbf{tidak
saling tumpang-tindih}. Sebaliknya, jika pada seluruh sumbu kandidat
proyeksinya selalu saling bertumpang-tindih, maka kedua objek tersebut
berpotongan.

Dalam konteks program ini, dua objek yang diuji adalah:
\begin{enumerate}
  \item sebuah \textbf{segitiga} yang merepresentasikan salah satu face
  dari model 3D, dan
  \item sebuah \textbf{AABB} yang merepresentasikan ruang voxel atau
  subruang octree yang sedang diperiksa.
\end{enumerate}

Tujuan pengujian ini adalah menentukan apakah suatu face masih relevan
untuk diteruskan ke suatu oktan. Jika segitiga tidak berpotongan dengan
AABB suatu oktan, maka face tersebut tidak perlu diproses lebih lanjut
pada cabang tersebut.

\subsection{Ide Dasar Pengujian}

Pengujian dilakukan dengan cara memproyeksikan segitiga dan AABB ke
sejumlah sumbu kandidat. Pada setiap sumbu, akan diperiksa apakah
interval proyeksi keduanya saling overlap. Apabila ditemukan satu saja
sumbu yang menghasilkan interval tidak overlap, maka sumbu tersebut
menjadi \textit{separating axis}, sehingga segitiga dan AABB dipastikan
tidak berpotongan.

Dengan kata lain:
\begin{itemize}
  \item jika ada sumbu pemisah, hasil interseksi adalah \texttt{false},
  \item jika tidak ada satu pun sumbu pemisah, hasil interseksi adalah
  \texttt{true}.
\end{itemize}

Karena baik segitiga maupun AABB merupakan bangun konveks, SAT dapat
diterapkan secara tepat dan efisien.

\subsection{Mengapa Menggunakan AABB?}

AABB adalah \textit{bounding box} yang sisinya selalu sejajar terhadap
sumbu koordinat X, Y, dan Z. Bentuk ini dipilih karena:
\begin{enumerate}
  \item representasinya sederhana, cukup dengan titik minimum dan
  maksimum,
  \item pembagian ruang menjadi delapan oktan dapat dilakukan dengan
  mudah,
  \item pengujian proyeksi terhadap sumbu X, Y, dan Z menjadi sangat
  efisien.
\end{enumerate}

Pada implementasi, AABB dinyatakan dengan dua titik:
\[
\texttt{min} = (x_{\min}, y_{\min}, z_{\min}), \qquad
\texttt{max} = (x_{\max}, y_{\max}, z_{\max})
\]
Dari dua titik ini dapat dihitung:
\begin{itemize}
  \item titik tengah AABB (\textit{center}),
  \item setengah ukuran box pada tiap sumbu (\textit{half extents}).
\end{itemize}

Segitiga kemudian ditranslasikan relatif terhadap pusat AABB agar
perhitungan proyeksi menjadi lebih sederhana.

\subsection{Tiga Belas Sumbu Kandidat}

Untuk pasangan segitiga dan AABB, terdapat \textbf{13 sumbu kandidat}
yang perlu diuji, yaitu:

\begin{enumerate}
  \item \textbf{3 sumbu normal face AABB}, yaitu sumbu X, Y, dan Z.

  Sumbu-sumbu ini digunakan untuk memeriksa apakah segitiga berada
  sepenuhnya di sebelah kiri/kanan, bawah/atas, atau depan/belakang box.
  Jika pada salah satu sumbu ini proyeksinya tidak overlap, maka
  interseksi langsung gagal.

  \item \textbf{1 sumbu normal bidang segitiga}.

  Sumbu ini diperoleh dari \textit{cross product} dua sisi segitiga.
  Pengujian pada sumbu ini memastikan apakah AABB dan bidang segitiga
  saling memotong atau berada pada sisi yang terpisah.

  \item \textbf{9 sumbu hasil cross product} antara 3 sumbu AABB dan 3
  sisi segitiga.

  Sumbu-sumbu ini diperlukan untuk menangani kasus ketika segitiga dan
  box tampak overlap pada sumbu utama X, Y, Z, tetapi sebenarnya tetap
  terpisah jika dilihat dari arah miring tertentu. Inilah bagian penting
  dari SAT yang membuat pengujian menjadi akurat untuk objek 3D.
\end{enumerate}


\subsection{Langkah Pengujian dalam Implementasi}

Secara umum, fungsi interseksi bekerja sebagai berikut:
\begin{enumerate}
  \item Hitung pusat AABB dan \textit{half extents}-nya.
  \item Translasi ketiga titik segitiga agar relatif terhadap pusat
  AABB.
  \item Hitung tiga vektor sisi segitiga.
  \item Uji 9 sumbu hasil \textit{cross product}.
  \item Uji 3 sumbu utama AABB, yaitu X, Y, dan Z.
  \item Uji 1 sumbu normal bidang segitiga.
  \item Jika semua pengujian lolos, maka segitiga dan AABB dinyatakan
  berpotongan.
\end{enumerate}

Setiap pengujian sumbu pada dasarnya membandingkan dua interval:
\begin{itemize}
  \item interval proyeksi tiga titik segitiga pada sumbu tersebut,
  \item interval proyeksi AABB pada sumbu yang sama.
\end{itemize}

Untuk segitiga, interval diperoleh dari nilai minimum dan maksimum hasil
\textit{dot product} titik-titik segitiga terhadap sumbu uji. Untuk
AABB, karena box berpusat di origin setelah translasi, jari-jari
proyeksinya dapat dihitung langsung dari kombinasi \textit{half extents}
dan nilai absolut komponen sumbu.

Jika interval segitiga berada sepenuhnya di luar interval AABB, maka
ditemukan sumbu pemisah dan fungsi segera mengembalikan
\texttt{false}.

Implementasi pada program ini mengikuti algoritma Möller (2001), dengan
optimasi kecil seperti penggunaan \texttt{std::min} dan \texttt{std::max}
bersarang untuk menghindari alokasi \texttt{initializer\_list} berulang.
Dengan demikian, uji interseksi tetap cukup cepat meskipun dieksekusi
dalam jumlah besar selama pembentukan octree.


\section{Analisis Kompleksitas}

\subsection{Kompleksitas Waktu}

\begin{itemize}
  \item \textbf{Kasus terbaik:} $O(F)$ — semua face berada di satu
        oktant, tidak ada subdivisi lanjutan.
  \item \textbf{Kasus rata-rata:} $O(F \cdot \log_8 N)$ — distribusi
        normal pada permukaan objek. Di sini $N = 8^n$ adalah jumlah
        voxel maksimal pada kedalaman $n$.
  \item \textbf{Kasus terburuk:} $O(F \cdot 8^n)$ — setiap face
        melintas semua oktant di setiap level. Tidak terjadi pada
        model nyata.
\end{itemize}

Tanpa optimasi, total node Octree adalah $\sum_{d=0}^{n} 8^d =
\frac{8^{n+1}-1}{7}$. Dengan \textit{pruning}, hanya node yang
bersinggungan dengan permukaan yang dikunjungi, sehingga jumlah
node aktual jauh lebih kecil.

\subsection{Kompleksitas Ruang}

$O(F + V)$ di mana $F$ adalah jumlah face aktif dalam rekursi saat ini
dan $V$ adalah jumlah voxel yang terbentuk. \textit{Stack} rekursi
memiliki kedalaman maksimal $n$ (parameter \texttt{maxDepth}).

\begin{table}[H]
  \centering
  \caption{Analisis Kompleksitas Waktu Algoritma}
  \label{tab:complexity}
  \begin{tabular}{lll}
    \toprule
    \textbf{Kasus} & \textbf{Kompleksitas Waktu} & \textbf{Keterangan} \\
    \midrule
    Terbaik  & $O(F)$                & Semua face di satu oktant \\
    Rata-rata & $O(F \cdot \log_8 N)$ & Distribusi normal permukaan \\
    Terburuk  & $O(F \cdot 8^n)$      & Setiap face melintas semua oktant \\
    \bottomrule
  \end{tabular}
\end{table}



% =============================================================
%  BAB 3 – IMPLEMENTASI
% =============================================================
\chapter{IMPLEMENTASI}

\section{Bahasa dan Tools}

\begin{table}[H]
  \centering
  \caption{Stack Teknologi yang Digunakan}
  \label{tab:tech}
  \begin{tabularx}{\textwidth}{lX}
    \toprule
    \textbf{Komponen} & \textbf{Detail} \\
    \midrule
    Bahasa       & C++17 \\
    Build System & CMake 3.14+ / Makefile \\
    Concurrency  & \texttt{std::async}, \texttt{std::future},
                   \texttt{std::mutex} (C++ Standard Library) \\
    Viewer       & OpenGL 1.x (Legacy), GLFW3 \\
    \bottomrule
  \end{tabularx}
\end{table}

\section{Struktur Program}

\subsection{Organisasi File}

\begin{table}[H]
  \centering
  \caption{Deskripsi Setiap File Source Code}
  \label{tab:files}
  \begin{tabularx}{\textwidth}{lX}
    \toprule
    \textbf{File} & \textbf{Deskripsi} \\
    \midrule
    \texttt{Vec3.hpp}
      & Struct vektor 3D dengan operasi dot, cross, normalize,
        dan operator overloading (\texttt{+}, \texttt{-}, \texttt{*}).
        Menggunakan \texttt{constexpr} dan \texttt{[[nodiscard]]}. \\[4pt]

    \texttt{AABB.hpp}
      & Struct \textit{Axis-Aligned Bounding Box}: fungsi
        \texttt{center()}, \texttt{halfExtents()},
        \texttt{intersectsTriangle()} (SAT Möller), dan
        \texttt{subdivide()} ke 8 oktant. \\[4pt]

    \texttt{ObjModel.hpp}
      & Container sederhana untuk daftar \textit{vertex} dan
        \textit{face} hasil parsing. \\[4pt]

    \texttt{ObjParser.hpp}
      & Namespace \texttt{ObjParser}: baca file \texttt{.obj},
        tangani format \texttt{v/vt/vn}, fan triangulation
        untuk poligon $>3$ vertex, validasi indeks. \\[4pt]

    \texttt{Octree.hpp}
      & Kelas utama: algoritma D\&C, \texttt{std::async}
        concurrency, pengumpulan voxel, dan statistik
        \textit{nodesFormed}/\textit{nodesSkipped}. \\[4pt]

    \texttt{ObjWriter.hpp}
      & Namespace \texttt{ObjWriter}: tulis voxel (AABB) ke
        file \texttt{.obj} — 8 vertex + 12 segitiga per voxel. \\[4pt]

    \texttt{Viewer.hpp}
      & Viewer 3D OpenGL/GLFW: orbit camera, render solid
        + wireframe, toggle model asli (efek X-Ray),
        implementasi manual \textit{LookAt} tanpa GLU. \\[4pt]

    \texttt{main.cpp}
      & Entry point: parsing argumen CLI,
        orkestrasi pipeline, output statistik. \\
    \bottomrule
  \end{tabularx}
\end{table}

\section{Modul \texttt{Vec3.hpp}}

Struct \texttt{Vec3} menyimpan koordinat \texttt{(x, y, z)} bertipe
\texttt{double} dan menyediakan operasi-operasi dasar vektor 3D yang
dibutuhkan oleh AABB dan renderer:

\begin{itemize}
  \item \texttt{operator+}, \texttt{operator-}, \texttt{operator*}
        — operasi aritmetika per-komponen.
  \item \texttt{dot()} — hasil kali titik, digunakan pada uji SAT.
  \item \texttt{cross()} — hasil kali silang, digunakan untuk menghitung
        normal bidang segitiga dan sumbu LookAt.
  \item \texttt{normalize()} — normalisasi dengan pengecekan panjang $> 0$.
  \item \texttt{operator<<} — overload untuk output stream langsung.
\end{itemize}

Seluruh operasi aritmetika dideklarasikan \texttt{constexpr} sehingga dapat
dievaluasi pada \textit{compile time} jika operand diketahui.

\section{Modul \texttt{AABB.hpp}}

\subsection{Fungsi \texttt{intersectsTriangle}}

Implementasi SAT sesuai Möller (2001) yang sudah dioptimasi.
Menggunakan \texttt{std::min}/\texttt{std::max} bersarang (bukan
\texttt{initializer\_list}) untuk menghindari alokasi heap berulang pada
loop rekursif yang intensif.

\subsection{Fungsi \texttt{subdivide}}

Mengembalikan \texttt{std::array<AABB, 8>} (bukan \texttt{vector}) sehingga
8 oktant anak dialokasikan di \textit{stack} tanpa overhead heap.

\section{Modul \texttt{Octree.hpp}}

Kelas \texttt{Octree} adalah inti program. Atribut utamanya:

\begin{itemize}
  \item \texttt{model} — referensi ke model yang di-parse.
  \item \texttt{maxDepth} — kedalaman rekursi maksimum.
  \item \texttt{voxels} — vektor AABB hasil akhir.
  \item \texttt{stats} — struct berisi \texttt{nodesFormed[]},
        \texttt{nodesSkipped[]}, dan \texttt{voxelCount}.
  \item \texttt{voxelMutex}, \texttt{statsMutex} — mutex untuk
        \textit{thread-safe access} saat concurrency aktif.
\end{itemize}

\subsection{Fungsi \texttt{build()}}

Menghitung bounding box seluruh model, membuat root AABB menjadi kubus
sempurna, mengisi \texttt{allFaces} menggunakan \texttt{std::iota},
dan memulai rekursi.

\subsection{Fungsi \texttt{buildRecursive()}}

Implementasi inti D\&C yang dijelaskan pada pseudocode di Bab 2.
Kunci desainnya adalah penggunaan \texttt{std::async} dengan
\texttt{std::launch::async} hanya ketika \texttt{depth < PARALLEL\_DEPTH\_LIMIT}
(= 2). Pada dua level teratas ini, hingga $8^2 = 64$ thread dapat
berjalan paralel, memberikan speedup nyata pada mesin multi-core.

\begin{lstlisting}[style=cpp, caption={Potongan implementasi concurrency},
  label={lst:async}]
if (currentDepth < PARALLEL_DEPTH_LIMIT) {
    std::array<std::future<void>, 8> threads;
    for (int i = 0; i < 8; ++i) {
        threads[i] = std::async(std::launch::async,
            [this, &subBoxes, &facesPerOktant, currentDepth, i]() {
                buildRecursive(subBoxes[i],
                               facesPerOktant[i],
                               currentDepth + 1);
            });
    }
    // Barrier: tunggu semua thread selesai
    for (auto& t : threads) t.get();
} else {
    for (int i = 0; i < 8; ++i)
        buildRecursive(subBoxes[i], facesPerOktant[i], currentDepth + 1);
}
\end{lstlisting}

\section{Modul \texttt{Viewer.hpp} (Bonus)}

Viewer dikompilasi secara kondisional dengan flag
\texttt{-DENABLE\_VIEWER}. Desain utamanya:

\begin{itemize}
  \item \textbf{Orbit camera:} State kamera disimpan dalam struct
        \texttt{ViewerState} yang diakses melalui pointer global
        (diperlukan karena callback GLFW bertipe function pointer C).
  \item \textbf{Manual LookAt:} Fungsi \texttt{applyLookAtMatrix()}
        diimplementasikan tanpa bergantung pada GLU agar kompatibel
        dengan Windows/MinGW yang sering tidak menyertakan GLU.
  \item \textbf{Efek X-Ray:} Saat toggle T aktif, model asli digambar
        dua kali — sekali sebelum voxel (tertimbun) dan sekali sesudah
        dengan \texttt{GL\_DEPTH\_TEST} dimatikan — menghasilkan efek
        \textit{tembus pandang}.
  \item \textbf{Reset otomatis:} \texttt{g\_state = ViewerState\{\}}
        mereset kamera setiap kali viewer dibuka.
\end{itemize}

\section{Modul \texttt{main.cpp}}

Entry point memisahkan logika parsing argumen ke fungsi
\texttt{parseArguments()} yang mengembalikan struct \texttt{AppConfig},
sehingga fungsi \texttt{main()} hanya berperan sebagai orkestrator
pipeline yang mudah dibaca.

\section{Cara Kompilasi dan Menjalankan Program}

\subsection{Kompilasi}

\begin{lstlisting}[style=pseudocode, numbers=none,
  caption={Cara kompilasi program}]
# Linux/Mac -- Tanpa viewer
mkdir -p bin
g++ -std=c++17 -O2 -o bin/voxelizer src/main.cpp -lpthread

# Linux -- Dengan viewer (GLFW harus terinstal)
g++ -std=c++17 -O2 -DENABLE_VIEWER -o bin/voxelizer_viewer \
    src/main.cpp -lpthread -lGL -lGLU -lglfw

# Windows (MSYS2/MinGW) -- Tanpa viewer
g++ -std=c++17 -O2 -o bin/voxelizer.exe src/main.cpp -lpthread
\end{lstlisting}

\subsection{Menjalankan Program}

\begin{lstlisting}[style=pseudocode, numbers=none,
  caption={Cara menjalankan program}]
# Tanpa viewer
./bin/voxelizer <path_file.obj> <max_depth>

# Dengan viewer
./bin/voxelizer_viewer <path_file.obj> <max_depth> --view

# Contoh
./bin/voxelizer test/input/cube.obj 3
./bin/voxelizer_viewer test/input/pumpkin.obj 6 --view
\end{lstlisting}

\subsection{Kontrol Viewer}

\begin{table}[H]
  \centering
  \caption{Kontrol Viewer 3D Interaktif}
  \label{tab:viewer}
  \begin{tabular}{ll}
    \toprule
    \textbf{Input} & \textbf{Aksi} \\
    \midrule
    Drag kiri mouse & Rotasi kamera (orbit) \\
    Scroll         & Zoom in/out \\
    Tombol T        & Toggle tampilkan model asli (efek X-Ray) \\
    Q / ESC         & Keluar viewer \\
    \bottomrule
  \end{tabular}
\end{table}

\section{Penjelasan Implementasi}

Implementasi utama program terletak pada kelas \texttt{Octree},
khususnya pada method \texttt{build()} dan \texttt{buildRecursive()}.
Kedua fungsi ini berperan dalam membangun struktur octree dan menghasilkan
voxel dari model 3D masukan.

Proses dimulai pada method \texttt{build()}, yaitu dengan menghitung
\textit{bounding box} yang mencakup seluruh vertex model. Selanjutnya,
\textit{bounding box} tersebut diubah menjadi kubus sempurna agar seluruh
voxel yang dihasilkan memiliki ukuran yang seragam. Setelah itu, program
menyusun daftar seluruh face pada model dan memulai proses rekursif dari
node akar dengan memanggil \texttt{buildRecursive()}.

Method \texttt{buildRecursive()} merepresentasikan inti algoritma
\textit{Divide and Conquer}. Pada setiap pemanggilan, fungsi menerima
sebuah AABB, himpunan face aktif yang masih berpotongan dengan area
tersebut, serta kedalaman saat ini. Jika himpunan face kosong, maka
cabang tersebut langsung dihentikan karena tidak mengandung bagian
permukaan objek. Jika kedalaman maksimum telah tercapai, AABB tersebut
langsung dicatat sebagai sebuah voxel.

Jika kedua kondisi tersebut belum terpenuhi, maka AABB saat ini dibagi
menjadi 8 oktan menggunakan fungsi \texttt{subdivide()}. Selanjutnya,
setiap face aktif diperiksa terhadap masing-masing oktan anak. Face yang
berpotongan dengan suatu oktan akan diteruskan ke node anak tersebut.
Dengan cara ini, setiap submasalah hanya memproses face yang relevan
dengan wilayah ruangnya.

Proses rekursif ini terus berlanjut hingga seluruh cabang mencapai kondisi
berhenti. Hasil akhirnya diperoleh dengan mengumpulkan seluruh
\textit{leaf node} yang valid ke dalam daftar \texttt{voxels}. Daftar
inilah yang kemudian menjadi representasi voxel dari model 3D awal.

Secara keseluruhan, implementasi ini memanfaatkan struktur octree untuk
membagi ruang secara bertahap dan hanya memproses bagian ruang yang
berkaitan dengan permukaan objek, sehingga proses voxelisasi menjadi lebih
efisien.

\section{Penjelasan Implementasi Spesifikasi Bonus}

\subsection{Bonus Concurrency}

Implementasi \textit{concurrency} terdapat pada kelas \texttt{Octree},
tepatnya di dalam method \texttt{buildRecursive()}. Pada proses
pembentukan octree, setiap node yang belum mencapai kondisi basis akan
dibagi menjadi 8 oktan. Delapan oktan ini merupakan submasalah yang
bersifat independen, sehingga dapat diproses secara bersamaan.

Pada implementasinya, ketika kedalaman rekursi masih lebih kecil dari
\texttt{PARALLEL\_DEPTH\_LIMIT}, pemrosesan kedelapan anak dijalankan
secara paralel menggunakan \texttt{std::async}. Setiap pemanggilan
\texttt{buildRecursive()} untuk satu oktan anak dibungkus dalam sebuah
\textit{task} asynchronous, lalu hasilnya disimpan dalam
\texttt{std::future}. Setelah seluruh \textit{task} dijalankan, program
menunggu semuanya selesai dengan memanggil \texttt{get()} pada setiap
\textit{future}.

Pendekatan ini dipilih karena pada level atas octree jumlah face yang
masih aktif umumnya besar, sehingga pemrosesan paralel dapat mempercepat
waktu eksekusi. Namun, paralelisme tidak diterapkan pada seluruh level.
Program membatasi concurrency hanya sampai kedalaman tertentu agar tidak
terjadi \textit{thread explosion} dan overhead yang terlalu besar.
Setelah melewati batas tersebut, rekursi dilanjutkan secara sekuensial.

Karena daftar voxel dan statistik dapat diakses oleh beberapa
\textit{thread} secara bersamaan, implementasi juga menggunakan
\texttt{mutex} untuk menjaga konsistensi data ketika menambahkan voxel
ke container hasil maupun saat memperbarui statistik node.

\subsection{Bonus Viewer 3D}

Implementasi \textit{viewer} 3D terdapat pada file \texttt{Viewer.hpp}
dan diaktifkan melalui target \texttt{voxelizer\_viewer}. Setelah proses
voxelisasi selesai, program dapat menampilkan model hasil dalam jendela
interaktif apabila dijalankan dengan flag \texttt{--view}.

Viewer dibangun menggunakan OpenGL dan GLFW. Data utama yang ditampilkan
adalah kumpulan voxel hasil konversi, yaitu daftar \texttt{AABB} yang
dihasilkan oleh octree. Setiap voxel divisualisasikan sebagai kubus 3D
dengan memanfaatkan fungsi \texttt{drawCube()}, sehingga pengguna dapat
melihat bentuk model hasil voxelisasi secara langsung.

Selain voxel, viewer juga dapat menampilkan model asli secara transparan
sebagai pembanding. Fitur ini membantu pengguna melihat seberapa baik
hasil voxelisasi merepresentasikan bentuk objek awal. Pada implementasi,
viewer menerima baik daftar voxel maupun data model asli, lalu
merender keduanya dalam satu scene 3D.

Interaksi pengguna juga disediakan agar tampilan tidak bersifat statis.
Kamera dapat diputar dengan mekanisme \textit{orbit}, zoom dapat diatur
menggunakan \textit{scroll mouse}, dan beberapa mode tampilan dapat
diubah melalui input keyboard, seperti menampilkan atau menyembunyikan
model asli serta mengganti mode \textit{wireframe}. Dengan demikian,
hasil voxelisasi tidak hanya disimpan ke file \texttt{.obj}, tetapi juga
dapat diamati secara visual dan interaktif langsung dari program.




% =============================================================
%  BAB 4 – EKSPERIMEN
% =============================================================
\chapter{EKSPERIMEN}

% ── Pengujian 1 ───────────────────────────────────────────────
\subsection{Test Case 1: \texttt{cow.obj} (depth = 4)}

\begin{figure}[H]
  \centering
  \includegraphics[width=0.85\textwidth]{cowcli.png}
  % \fbox{\parbox{0.85\textwidth}{\centering\vspace{2.5cm}
  %   [Screenshot terminal: perintah + output statistik cube depth 2]
  %   \vspace{2.5cm}}}
  \caption{Input dan output CLI untuk cow.obj (depth = 4)}
\end{figure}

\begin{figure}[H]
  \centering
  \includegraphics[width=0.7\textwidth]{cow3d.png}
  % \fbox{\parbox{0.7\textwidth}{\centering\vspace{3cm}
  %   [Screenshot hasil voxel cube.obj di viewer / OBJ viewer]
  %   \vspace{3cm}}}
  \caption{Hasil voxelization cow.obj (depth = 4)}
\end{figure}

% ── Pengujian 2 ───────────────────────────────────────────────
\subsection{Test Case 2: \texttt{house.obj} (depth = 6)}

\begin{figure}[H]
    \centering
  \includegraphics[width=0.85\textwidth]{housecli.png}
  % \fbox{\parbox{0.85\textwidth}{\centering\vspace{2.5cm}
  %   [Screenshot terminal: perintah + output statistik cube depth 2]
  %   \vspace{2.5cm}}}
  \caption{Input dan output CLI untuk house.obj (depth = 6)}
\end{figure}

\begin{figure}[H]
  \centering
    \includegraphics[width=0.7\textwidth]{house3d.png}
  % \fbox{\parbox{0.85\textwidth}{\centering\vspace{2.5cm}
  %   [Screenshot terminal + hasil voxel house.obj depth 4]
  %   \vspace{2.5cm}}}
  \caption{Hasil voxelization house.obj (house = 6)}
\end{figure}

% ── Pengujian 3 ───────────────────────────────────────────────
\subsection{Test Case 3: \texttt{line.obj} (depth = 3)}

\begin{figure}[H]
    \centering
  \includegraphics[width=0.85\textwidth]{linecli'.png}
  % \fbox{\parbox{0.85\textwidth}{\centering\vspace{2.5cm}
  %   [Screenshot terminal: perintah + output statistik cube depth 2]
  %   \vspace{2.5cm}}}
  \caption{Input dan output CLI untuk line.obj (depth = 3)}
\end{figure}

\begin{figure}[H]
  \centering
    \includegraphics[width=0.7\textwidth]{line3d.png}
  % \fbox{\parbox{0.85\textwidth}{\centering\vspace{2.5cm}
  %   [Screenshot terminal + hasil voxel house.obj depth 4]
  %   \vspace{2.5cm}}}
  \caption{Hasil voxelization line.obj (depth = 3)}
\end{figure}

% ── Pengujian 5 ───────────────────────────────────────────────
\subsection{Test Case 4: \texttt{pumpkin.obj} (depth = 5)}

\begin{figure}[H]
    \centering
  \includegraphics[width=0.85\textwidth]{pumpcli.png}
  % \fbox{\parbox{0.85\textwidth}{\centering\vspace{2.5cm}
  %   [Screenshot terminal: perintah + output statistik cube depth 2]
  %   \vspace{2.5cm}}}
  \caption{Input dan output CLI untuk pumpkin.obj (depth = 5)}
\end{figure}

\begin{figure}[H]
  \centering
    \includegraphics[width=0.7\textwidth]{pump3d.png}
  % \fbox{\parbox{0.85\textwidth}{\centering\vspace{2.5cm}
  %   [Screenshot terminal + hasil voxel house.obj depth 4]
  %   \vspace{2.5cm}}}
  \caption{Hasil voxelization pumpkin.obj (line = 5)}
\end{figure}

\subsection{Test Case 5: \texttt{teapot.obj} (depth = 7)}

\begin{figure}[H]
    \centering
  \includegraphics[width=0.85\textwidth]{teapotcli.png}
  % \fbox{\parbox{0.85\textwidth}{\centering\vspace{2.5cm}
  %   [Screenshot terminal: perintah + output statistik cube depth 2]
  %   \vspace{2.5cm}}}
  \caption{Input dan output CLI untuk teapot.obj (depth = 7)}
\end{figure}

\begin{figure}[H]
  \centering
    \includegraphics[width=0.7\textwidth]{teapot3d.png}
  % \fbox{\parbox{0.85\textwidth}{\centering\vspace{2.5cm}
  %   [Screenshot terminal + hasil voxel house.obj depth 4]
  %   \vspace{2.5cm}}}
  \caption{Hasil voxelization teapot.obj (depth = 7)}
\end{figure}


% =============================================================
%  BAB 5 – KESIMPULAN DAN SARAN
% =============================================================
\chapter{Penutup}

\section{Kesimpulan}

Dari pengerjaan Tugas Kecil 2 ini, dapat diambil beberapa kesimpulan yaitu, program voxelization 3D berhasil diimplementasikan menggunakan algoritma
Divide and Conquer berbasis Octree dalam bahasa C++17. Algoritma bekerja
dengan membagi ruang secara rekursif menjadi 8 oktant, memangkas oktant
yang tidak berisi permukaan (\textit{pruning}), dan menghasilkan voxel pada
leaf node yang bersinggungan dengan segitiga model. Implementasi terbukti
menghasilkan output yang benar dan seragam (\textit{uniform voxels}).

Pruning berdasarkan hasil uji SAT secara signifikan mengurangi jumlah node
yang dikunjungi dibandingkan traversal penuh. Pada model dengan permukaan
tipis (\textit{thin shell}), hanya sebagian kecil dari $8^n$ node maksimal
yang benar-benar diproses.

Selain itu Penggunaan \texttt{std::async} pada dua level kedalaman teratas Octree
memberikan speedup nyata pada mesin multi-core tanpa memerlukan library
eksternal. Batas \texttt{PARALLEL\_DEPTH\_LIMIT = 2} mencegah
\textit{thread explosion} di kedalaman bawah.

Viewer interaktif berbasis OpenGL/GLFW memudahkan verifikasi visual hasil
voxelization. Implementasi manual \textit{LookAt} tanpa GLU meningkatkan
portabilitas di berbagai platform.


\section{Saran}
Pelaksanaan Tugas Kecil 2 IF2211 Strategi Algoritma di Semester II (Genap)
Tahun 2024/2025 merupakan pengalaman yang sangat berharga bagi kami. Dari
pengalaman ini, kami ingin berbagi beberapa saran kepada pembaca yang mungkin
akan menghadapi tugas serupa di masa depan:
\begin{infobox}
\textbf{Ucup:}
Terimakasih asisten yang sudah memberi tucil di tengah libur lebaran.
\end{infobox}

\begin{infobox}
\textbf{Alvan:} Saya pikir tubes ini didesain kepada kita untuk melatih
diri menghadapi tekanan deadline yang mungkin nanti ketika di industri
akan relate. Namun, tugas ini sebaiknya dikerjakan sesegera mungkin.
Dengan waktu kurang lebih 2 minggu saya merasa cukup chaos menghadapi
tugas-tugas yang ada. Rasa-rasanya setelah tugas ini selesai pun
ketenangan tidak serta merta menghampiri, terpaan deadline selalu
menghantui. Pagi, siang malam, dilanda ketakutan akan tugas. Bingung
memikir bagaimana ku menghadapinya. Namun, aku memiliki solusi dan saran
kepadamu, ``Sebaiknya kalo ada tugas minimal di-extend 5 hari''.
\end{infobox}

% =============================================================
%  LAMPIRAN
% =============================================================
\section{Tautan}

Seluruh implementasi kode program tersedia di repository GitHub:

\begin{infobox}
\textbf{Repository:}
\texttt{https://github.com/ucupsss/Tucil2\_13524014\_13524124}

\end{infobox}

\section{Lampiran}

Berikut adalah tabel checklist pemenuhan spesifikasi tugas:

\vspace{0.5cm}
\begin{table}[H]
  \centering
  \begin{tabular}{clcc}
    \toprule
    \textbf{No} & \textbf{Poin} & \textbf{Ya} & \textbf{Tidak} \\
    \midrule
    1 & Program berhasil dikompilasi tanpa kesalahan
      & \cmark & \\
    2 & Program berhasil dijalankan
      & \cmark & \\
    3 & Hasil konversi \texttt{.obj} terdiri dari kubus-kubus kecil
      & \cmark & \\
    4 & Program menampilkan informasi/\textit{report} di CLI
      & \cmark & \\
    5 & Program menerima masukan \texttt{.obj} dan mengembalikan output
      & \cmark & \\
    6 & \textbf{[Bonus]} Program menggunakan \textit{concurrency}
      & \cmark & \\
    7 & \textbf{[Bonus]} Program memiliki \texttt{.obj} viewer
      & \cmark & \\
    \bottomrule
  \end{tabular}
  \caption{Checklist kelengkapan program}
  \label{tab:checklist}
\end{table}

\chapter{Pernyataan tidak melakukan kecurangan yang ditandatangani}

\vspace{1cm}
\begin{mdframed}[linewidth=1pt]
  \begin{center}
    \textbf{PERNYATAAN}
  \end{center}

  \vspace{0.5cm}
  Kami yang bertanda tangan di bawah ini menyatakan bahwa Tugas ini disusun sepenuhnya tanpa bantuan kecerdasan buatan (Generative AI), melainkan hasil pemikiran dan analisis mandiri.


  \vspace{2cm}
  \begin{flushright}
    \begin{figure}[H]
        \raggedleft
      \includegraphics[width=0.3\textwidth]{ttducup.png}
    \end{figure}
    \textbf{Yusuf Faishal Listyardi} \\
        \begin{figure}[H]
        \raggedleft
      \includegraphics[width=0.3\textwidth]{WhatsApp Image 2026-03-27 at 23.16.56.jpeg}
    \end{figure}
    \textbf{Zahran Alvan Putra Winarko}
  \end{flushright}
\end{mdframed}

\end{document}
